%% Methods in Ecology and Evolution — Application Paper
%% Run LaTeX on this file several times to get cross-references and citations.

\documentclass[11pt]{book}
\usepackage{Wiley-AuthoringTemplate}
\usepackage[authoryear]{natbib}

%%--------------------------------------------------------------------
%% Section numbering depth
\setcounter{secnumdepth}{3}

%%--------------------------------------------------------------------
%% Additional packages not in the Wiley template
\usepackage{amsmath,amssymb}
\usepackage{booktabs}
\usepackage{listings}
\usepackage{xcolor}
\usepackage{caption}
\usepackage{microtype}
\hypersetup{colorlinks=true,linkcolor=blue!60!black,citecolor=blue!60!black,urlcolor=blue!60!black}

%%--------------------------------------------------------------------
%% Code listing style
\lstset{
  basicstyle=\ttfamily\small,
  backgroundcolor=\color{gray!8},
  frame=single,
  framerule=0.4pt,
  rulecolor=\color{gray!40},
  breaklines=true,
  breakatwhitespace=true,
  columns=fullflexible,
  keepspaces=true,
  showstringspaces=false,
  tabsize=4,
  captionpos=b,
  numbers=none,
}
\lstdefinestyle{python}{language=Python}
\lstdefinestyle{bash}{language=bash}

% =============================================================================
\begin{document}

\mainmatter

% =============================================================================
\chapter[\textbf{gdmbayes}: A Python Package for Bayesian Generalised
         Dissimilarity Modelling]%
        {\textbf{gdmbayes}: A Python Package for Bayesian Generalised
         Dissimilarity Modelling}

\author*[1]{Harold Horsley}

\address[1]{\orgdiv{[Department]},
\orgname{[Institution]},
\postcode{[Postcode]},
\city{[City]},
\country{[Country]}}

\address*{Corresponding Author: Harold Horsley; \email{[email]}}

\maketitle

% ---- Abstract ---------------------------------------------------------------
\begin{abstract}{Abstract}
Generalised dissimilarity modelling (GDM) is a widely used framework for
relating pairwise ecological dissimilarities between sites to environmental
and spatial predictors using non-linear I-spline basis functions. Despite its
prevalence in ecology, no Python GDM implementation has existed and the method
has been confined to a single R package providing only point estimates with no
uncertainty quantification. Propagating ecological uncertainty through predictor
importance or incorporating spatial random effects has therefore not been
possible within any existing GDM tooling.

We present \textbf{gdmbayes}, an open-source Python package providing both a
frequentist scikit-learn-compatible GDM estimator and a full Bayesian extension
built on PyMC. The Bayesian estimator supports flexible variance structures,
Gaussian process spatial random effects, and hierarchical predictor importance
weights within a reproducible, configuration-driven workflow. We demonstrate
the package on the canonical southwest Australia plant-diversity dataset,
replicating and extending the benchmarks of \citet{white2024} and show that
the frequentist estimator achieves a 10-fold cross-validated RMSE of 0.064
compared with 0.074 for the R \textbf{gdm} implementation.

\textbf{gdmbayes} is the first Python GDM implementation and the only
packaged software---in any language---for full Bayesian GDM with spatial
random effects. The package is freely available at
\url{https://github.com/harryhorsley9/spgdmm} under the Apache~2.0 licence.
\end{abstract}

\keywords{Bayesian inference; beta diversity; compositional turnover; ecological
  dissimilarity; generalised dissimilarity modelling; Gaussian process;
  Python; scikit-learn; species distribution modelling}

% ---- 1. Introduction --------------------------------------------------------
\section{Introduction}

Ecologists increasingly rely on large spatially referenced datasets to
understand how species composition changes across landscapes. Generalised
dissimilarity modelling \citep[GDM;][]{ferrier2007} has become one of the most
widely applied frameworks for this purpose. GDM relates pairwise ecological
dissimilarities---typically Bray--Curtis or S{\o}rensen distances computed from
species assemblage data---to pairwise differences in environmental predictors
and geographic distance using a generalised additive model built on I-spline
basis functions. The non-linear I-spline representation allows GDM to capture
the empirical observation that the rate of compositional turnover saturates at
high environmental distances, a biological property that linear models cannot
accommodate \citep{ferrier2007,mokany2022}.

GDM has been applied to beta-diversity gradients, conservation planning
\citep{mokany2022}, climate-change vulnerability assessment
\citep{fitzpatrick2013}, and spatial prioritisation of biodiversity surveys.
The R package \textbf{gdm} \citep{fitzpatrick2021} has been the reference
implementation for nearly two decades, providing a point-estimate frequentist
estimator via non-negative least squares (NNLS). Two important limitations have
persisted. First, no Python implementation exists, which limits integration
with the scientific Python ecosystem and with machine-learning pipelines built
around scikit-learn. Second, neither the R package nor any other tool provides
uncertainty estimates for GDM coefficients or predictions, making it impossible
to propagate ecological uncertainty through downstream analyses.

Several Bayesian extensions of GDM have been proposed. \citet{woolley2017}
introduced BBGDM, which replaces NNLS with Bayesian bootstrap resampling; this
provides interval estimates but is not proper Bayesian inference (no likelihood
or priors). \citet{lusiana2023} derived a fully Bayesian GDM using
Hamiltonian Monte Carlo in \textbf{rstan} with a binomial likelihood and logit
link, but provide no reusable software package. Most recently, \citet{white2024}
introduced the spatial GDM mixed model (spGDMM), a hierarchical Bayesian
formulation with a censored-normal (Tobit) likelihood, Gaussian process spatial
random effects, and 10-fold cross-validation model selection, implemented in
NIMBLE for R. This provides the strongest statistical foundation to date but
remains inaccessible to Python users.

Here we present \textbf{gdmbayes}, a Python package that provides (i)~a
frequentist GDM estimator replicating the R \textbf{gdm} algorithm and fully
compatible with the scikit-learn API \citep{pedregosa2011}, and (ii)~a
Bayesian GDM estimator implementing the spGDMM model of \citet{white2024}
using PyMC \citep{salvatier2016} and ArviZ \citep{kumar2019} for MCMC
sampling and posterior diagnostics. Both estimators share a common
preprocessing layer that handles I-spline mesh construction, pairwise distance
computation, and feature matrix assembly.

% ---- 2. Package Description -------------------------------------------------
\section{Package Description}

\subsection{Overview and Installation}

\textbf{gdmbayes} is organised around three main classes---\texttt{GDM},
\texttt{spGDMM}, and \texttt{GDMModel}---sharing a common preprocessing
component (\texttt{GDMPreprocessor}). All classes follow the scikit-learn
\texttt{fit(X,~y)} / \texttt{predict(X)} API.

Installation via pip:

\begin{lstlisting}[style=bash]
pip install gdmbayes
\end{lstlisting}

Dependencies include PyMC~$\geq$~5.0, ArviZ, nutpie, NumPy, pandas, SciPy,
xarray, and scikit-learn. A conda environment file is provided for
reproducibility.

\subsection{Preprocessing}

All preprocessing is handled by \texttt{GDMPreprocessor}, a standalone
scikit-learn-compatible transformer (\texttt{BaseEstimator},
\texttt{TransformerMixin}). The preprocessor computes the I-spline feature
matrix from raw site-level data and stores fitted state (knot meshes, location
arrays) for use at prediction time.

\paragraph{I-spline basis functions.}
For each environmental predictor $x$, the preprocessor constructs
$J = \texttt{deg} + \texttt{knots}$ I-spline basis functions
$I_1(x), \ldots, I_J(x)$ over a mesh of evenly spaced or percentile-based knot
positions. I-splines are the integral of M-splines and are monotone
non-decreasing by construction \citep{ramsay1988}, ensuring predicted
dissimilarity increases monotonically with environmental distance. The pairwise
feature for predictor $k$ between sites $i$ and $j$ is:

\begin{equation}
  f_{k,\ell}(i,j) = \bigl|I_\ell(x_{k,i}) - I_\ell(x_{k,j})\bigr|,
  \quad \ell = 1, \ldots, J.
  \label{eq:ispline-feature}
\end{equation}

Geographic distance is optionally included as an additional predictor using the
same I-spline transformation applied to pairwise Euclidean, geodesic, or
grid-based ocean distances. Preprocessing is configured via
\texttt{PreprocessorConfig}:

\begin{lstlisting}[style=python]
from gdmbayes import PreprocessorConfig

cfg = PreprocessorConfig(
    deg=3,                       # I-spline polynomial degree
    knots=2,                     # number of internal knots
    mesh_choice="percentile",    # "percentile", "even", or "custom"
    distance_measure="geodesic", # "euclidean", "geodesic", "ocean_distance"
)
\end{lstlisting}

\subsection{Frequentist GDM}
\label{sec:frequentist}

\texttt{GDM} implements the standard GDM algorithm from \citet{ferrier2007} as
a scikit-learn estimator. Given the pairwise I-spline feature matrix
$\mathbf{X}_{\mathrm{GDM}} \in \mathbb{R}^{P \times Q}$ ($P$ site pairs, $Q$
I-spline columns), fitting proceeds by iteratively reweighted NNLS (matching
the C implementation in the R \textbf{gdm} package):

\begin{equation}
  \min_{\boldsymbol{\beta} \geq 0}
  \bigl\| \mathbf{W}^{1/2}(g(\mathbf{y}) - \mathbf{X}_{\mathrm{GDM}}\boldsymbol{\beta}) \bigr\|^2,
  \label{eq:nnls}
\end{equation}

where $g(y) = -\log(1-y)$ is the complementary log-log link and $\mathbf{W}$
is the diagonal matrix of IRLS weights derived from the binomial variance
function $V(\mu) = \mu(1-\mu)$. Jeffreys smoothing initialises $\eta$ to
avoid divergence when $y = 1$. Predicted dissimilarities are recovered via:

\begin{equation}
  \hat{y} = 1 - \exp\!\bigl(-\mathbf{X}_{\mathrm{GDM}}\hat{\boldsymbol{\beta}}\bigr).
  \label{eq:inverse-link}
\end{equation}

Predictor importance is the sum of NNLS coefficients over all I-spline basis
functions per predictor, analogous to the R \textbf{gdm} package.

\begin{lstlisting}[style=python]
from gdmbayes import GDM

model = GDM(geo=True, splines=3, knots=2)
model.fit(X, y)
print(f"Deviance explained: {model.explained_:.4f}")
predictions = model.predict(X_new)
\end{lstlisting}

\subsection{Bayesian GDM (\texttt{spGDMM})}
\label{sec:bayesian}

\texttt{spGDMM} extends the frequentist GDM with full Bayesian inference via
PyMC, implementing the spGDMM model of \citet{white2024}. The response is
modelled on the log scale with upper censoring at $y = 1$ (the Tobit
formulation):

\begin{equation}
  \log(y_{ij}) \sim \mathrm{CensoredNormal}\!\left(\mu_{ij},\;
  \sigma^2_{ij},\; \text{upper} = 0\right).
  \label{eq:likelihood}
\end{equation}

\paragraph{Linear predictor.}
With the \texttt{alpha\_importance} option, the linear predictor for site pair
$(i, j)$ is:

\begin{equation}
  \mu_{ij} = \beta_0
  + \sum_{k=1}^{K} \alpha_k \sum_{\ell=1}^{J} \beta_{k,\ell}\, f_{k,\ell}(i,j)
  + \eta_{ij},
  \label{eq:linear-predictor}
\end{equation}

where $\beta_0 \sim \mathcal{N}(0, 10)$; the normalised I-spline coefficients
$\boldsymbol{\beta}_k \sim \mathrm{Dirichlet}(\mathbf{1})$ ensure
$\sum_\ell \beta_{k,\ell} = 1$; and per-predictor importance weights
$\alpha_k \sim \mathrm{HalfNormal}(1)$ capture total magnitude. The term
$\eta_{ij}$ is the optional spatial random effect (Section~\ref{sec:spatial}).

\paragraph{Spatial random effects.}
\label{sec:spatial}
An optional Gaussian process prior models spatial autocorrelation:

\begin{equation}
  \boldsymbol{\psi} \sim \mathcal{GP}\!\left(0,\;
  \sigma^2_\psi \cdot \mathrm{Exp}\!\left(\|\cdot\| / \ell\right)\right),
  \label{eq:gp-prior}
\end{equation}

placed on the $n$ training sites. The pairwise spatial contrast
$\eta_{ij} = |\psi_i - \psi_j|$ or $(\psi_i - \psi_j)^2$ is added to
$\mu_{ij}$. The length scale $\ell$ is set to half the median inter-site
distance; $\sigma^2_\psi \sim \mathrm{InverseGamma}(1,1)$.

\paragraph{Sampler.}
MCMC uses the No-U-Turn Sampler \citep[NUTS;][]{hoffman2014} via PyMC, with
nutpie as the default backend. Sampler settings are controlled via
\texttt{SamplerConfig}:

\begin{lstlisting}[style=python]
from gdmbayes import spGDMM, ModelConfig, PreprocessorConfig, SamplerConfig

model = spGDMM(
    preprocessor=PreprocessorConfig(deg=3, knots=2, distance_measure="geodesic"),
    model_config=ModelConfig(variance="homogeneous", spatial_effect="abs_diff"),
    sampler_config=SamplerConfig(draws=1000, tune=1000, chains=4, random_seed=42),
)
idata = model.fit(X, y)
\end{lstlisting}

Results are stored as an ArviZ \texttt{InferenceData} object, enabling the
full ArviZ diagnostic toolkit (R-hat, effective sample size, posterior
predictive checks). Fitted models are serialised as NetCDF files via
\texttt{model.save(fname)} for prediction on new data after deserialisation.

% ---- 3. Running Example -----------------------------------------------------
\section{Running Example: Southwest Australia Plant Diversity}

We demonstrate \textbf{gdmbayes} on the canonical southwest Australia dataset
from the R \textbf{gdm} package, the same dataset used in \citet{white2024}
Table~1. This dataset comprises 94 sites with occurrence data for 856 plant
species and 10 environmental predictors \citep{fitzpatrick2013}, with
Bray--Curtis dissimilarities ($n = 4{,}371$ site pairs). We compare against
published benchmarks from the Ferrier et al. (R \textbf{gdm}) and spGDMM
models in \citet{white2024} Table~1.

\subsection{Frequentist GDM}

\begin{lstlisting}[style=python]
import pandas as pd
from gdmbayes import GDM, PreprocessorConfig

sites = pd.read_csv("data/southwest_sites.csv")
y = pd.read_csv("data/southwest_y.csv")["y"].values

X = sites.rename(columns={"Long": "xc", "Lat": "yc"}).copy()
X["time_idx"] = 0
X = X[["xc", "yc", "time_idx"] + ENV_PREDICTORS]

model = GDM(
    geo=True,
    preprocessor_config=PreprocessorConfig(
        deg=3, knots=2, mesh_choice="percentile", distance_measure="geodesic"
    ),
)
model.fit(X, y)
print(f"Deviance explained: {model.explained_:.4f}")
\end{lstlisting}

Table~\ref{tab:benchmark} shows predictive performance under 10-fold
cross-validation on site pairs (the methodology of \citealt{white2024}):
fitting the preprocessor on all 94 sites and the NNLS step on training pairs
only, then predicting withheld pairs. \textbf{gdmbayes} \texttt{GDM} achieves
RMSE~=~0.064 and MAE~=~0.042, outperforming the R \textbf{gdm} benchmark
(RMSE~=~0.074, MAE~=~0.055; \citealt{white2024}).

\begin{table}[h]
\caption{Predictive performance on the southwest Australia dataset
  (10-fold cross-validation on site pairs), compared with benchmarks from
  \citet{white2024} Table~1. RMSE and MAE are on the Bray--Curtis dissimilarity
  scale ($y \in [0, 1]$).\label{tab:benchmark}}{%
\begin{tabular}{@{}lcc@{}}
\toprule
Model & RMSE & MAE \\
\midrule
R \textbf{gdm} --- Ferrier \citep{white2024} & 0.0737 & 0.0549 \\
Best spGDMM \citep{white2024}                 & 0.0731 & 0.0545 \\
\textbf{gdmbayes} \texttt{GDM} (this paper)  & 0.0636 & 0.0424 \\
\textbf{gdmbayes} \texttt{spGDMM} --- no spatial & [TBD] & [TBD] \\
\textbf{gdmbayes} \texttt{spGDMM} --- abs\_diff  & [TBD] & [TBD] \\
\botrule
\end{tabular}}{}
\end{table}

\subsection{Bayesian spGDMM}

\begin{lstlisting}[style=python]
from gdmbayes import spGDMM, ModelConfig, SamplerConfig, PreprocessorConfig

model = spGDMM(
    preprocessor=PreprocessorConfig(
        deg=3, knots=2, mesh_choice="percentile", distance_measure="geodesic"
    ),
    model_config=ModelConfig(
        variance="homogeneous", spatial_effect="abs_diff", alpha_importance=True
    ),
    sampler_config=SamplerConfig(
        draws=1000, tune=1000, chains=4, target_accept=0.95,
        nuts_sampler="nutpie", random_seed=42,
    ),
)
idata = model.fit(X, y)
model.save("southwest_bayes_absdiff.nc")
\end{lstlisting}

% ---- 4. Comparison ----------------------------------------------------------
\section{Comparison with Existing Tools}

Table~\ref{tab:comparison} compares \textbf{gdmbayes} with the R \textbf{gdm}
package \citep{fitzpatrick2021} and prior Bayesian GDM work.

\begin{table}[h]
\caption{Feature comparison of GDM implementations across key dimensions.
  \textbf{gdmbayes} is the only Python implementation and the only packaged
  software for the full Bayesian spGDMM model.\label{tab:comparison}}{%
\begin{tabular}{@{}lccccc@{}}
\toprule
Feature & R \textbf{gdm} & BBGDM & Lusiana & White spGDMM & \textbf{gdmbayes} \\
        &                & \citep{woolley2017} & \citep{lusiana2023} & \citep{white2024} & \\
\midrule
Language               & R   & R  & R     & R (NIMBLE) & \textbf{Python} \\
Frequentist GDM        & Yes & Yes & Yes  & Yes & Yes \\
Proper Bayesian inference & No & No & Yes & Yes & Yes \\
Spatial random effects & No  & No  & No   & Yes & Yes \\
One-inflation (Tobit)  & No  & No  & No   & Yes & Yes \\
sklearn API            & No  & No  & No   & No  & Yes \\
Installable package    & Yes & Yes & No   & No  & Yes \\
\botrule
\end{tabular}}{}
\end{table}

The R \textbf{gdm} package remains the reference for frequentist GDM and
provides additional utilities (partial regression plots, Monte Carlo
significance testing) not yet included in \textbf{gdmbayes}. BBGDM
\citep{woolley2017} provides bootstrap intervals but is not fully Bayesian.
The Lusiana BGDM \citep{lusiana2023} uses a simpler binomial--logit
specification with no spatial effects and no reusable software package.
\textbf{gdmbayes} implements the full spGDMM model of \citet{white2024} as the
first packaged, cross-platform Python software for hierarchical Bayesian GDM.

% ---- 5. Availability --------------------------------------------------------
\section{Availability}

\textbf{gdmbayes} is available at \url{https://github.com/harryhorsley9/spgdmm}
under the Apache~2.0 licence. A PyPI release is forthcoming. A Zenodo DOI will
be minted at the time of publication. The package requires Python~$\geq$~3.10
and is tested on Linux. The southwest Australia dataset used in the running
example is available from the R \textbf{gdm} package (\texttt{data(southwest)},
\citealt{fitzpatrick2021}).

% ---- Acknowledgements -------------------------------------------------------
\section*{Acknowledgements}

[To be completed.]

% ---- Author Contributions ---------------------------------------------------
\section*{Author Contributions}

H.~Horsley: conceptualisation, software, formal analysis, writing --
original draft, writing -- review and editing.

% ---- Data Availability ------------------------------------------------------
\section*{Data Availability}

The southwest Australia site and dissimilarity data used in the running example
are derived from the R \textbf{gdm} package \citep{fitzpatrick2021} and are
included in the \texttt{examples/data/} directory of the repository. All
example scripts (\texttt{examples/southwest\_example.py},
\texttt{examples/bci\_example.py}) are fully self-contained and reproducible.

% ---- References -------------------------------------------------------------
\backmatter
\bibliographystyle{plainnat}
\bibliography{references}

\end{document}
